\documentclass[10pt,a4paper]{amsart}
\usepackage[margin=19mm]{geometry}
\usepackage{amsmath,amssymb,mathtools}
\usepackage[scr=rsfs,cal=euler]{mathalfa}
\usepackage{xfrac}
\usepackage[unicode,hidelinks,bookmarks=false]{hyperref}
\newcommand{\ie}{i.e.\ }
\newcommand{\cf}{cf.\ }
\newcommand{\resp}{respectively\ }
\newcommand{\Subsection}{\subsection}
\providecommand{\nobreakdash}{\nobreak\hbox{-}}
\setlength{\emergencystretch}{2em}
\multlinegap0pt
\usepackage{amssymb}
\usepackage{xcolor}
\usepackage{algorithm}\usepackage{algpseudocode}
\newtheorem{defi}{Definition}
\newtheorem{thm}{Theorem}
\title{On linearly ordered sets of chain components}
\author{P. Cintioli, A. Della Corte, M. Farotti}
\date{Source: arXiv:2507.21798v3 (CC BY 4.0)}
\begin{document}
\maketitle

\noindent\emph{Source and licence.} On linearly ordered sets of chain components. Version arXiv:2507.21798v3 (CC BY 4.0), distributed by arXiv under the Creative Commons Attribution 4.0 International licence, http://creativecommons.org/licenses/by/4.0/, as recorded in the arXiv licence metadata for this exact version and shown on the abstract page https://arxiv.org/abs/2507.21798v3. Full licence notice: \texttt{licenses/arxiv-2507.21798-CC-BY-4.0.txt}.\par\smallskip

\noindent\textit{Editorial scope.} Counted unit: the article's thm order\_type\_[0,1] (source0.tex lines 523--525). The article's complete original proof of it is delivered with it (source0.tex lines 526--667, one \texttt{proof} environment of 142 lines). No mathematics is written, repaired or re-derived: the statement and the proof are the article's own lines, in the article's own notation, and no retained line is substituted or re-flowed.  Retained uncounted as the source-local context the proof needs: lines 517--521.  Nothing was omitted from any retained span, so the package carries no omission record.  No retained line contains a citation, so the wrapper carries no bibliography and no bibliography entry.  No retained line contains a figure environment, an image-inclusion command or drawing code, so no image is delivered and the package declares no asset dependency.  Editorial formatting only, in addition: an AMS \texttt{amsart} preamble that restates the article's own declarations and macros used by the retained lines, namely the environments \texttt{defi}, \texttt{thm} on separate counters, together with the standard packages those lines need.  The retained environments print in this document as Definition 1, Theorem 1 in order of appearance, whereas the article prints the same items with the numbering of its own full document.  The complete native source of the article as distributed in the arXiv e-print for this exact version is bundled unchanged as \texttt{sources/182365/source0.tex}.
\begin{defi}\label{def_order_power}
    Let $\mathcal{P}([0,1])$ be the power set of the unit interval. 
    Let us define a partial order relation on $\mathcal{P}([0,1])$, denoted by $\le_{\mathcal{P}}$, as follows.
    For $I,J\in \mathcal{P}([0,1])$, we say that $I\leq_{\mathcal{P}} J$ if and only if either $I=J$ or $x<y$ for every $x\in I$ and $y\in J$. We say that $I <_\mathcal{P} J$ if $I\le_\mathcal{P} J$ and $I\neq J$.
\end{defi}

\begin{thm}\label{order_type_[0,1]}
There is a discontinuous dynamical system $([0,1],f)$ such that the poset $(\mathfrak{C}_f,\preceq)$ has order type $([0,1]\cap\mathbb{Q},\le)$.
\end{thm}

\begin{proof}
We want to define a dynamical system $([0,1],f)$ such that its chain components poset is order isomorphic to $(L,\leq_{\mathcal{P}})$, which is in turn order isomorphic to $([0,1]\cap \mathbb{Q}, \leq)$, where $L=\{[a_n,b_n]\}_{n\geq 0}$ is a countable set of pairwise disjoint intervals in $[0,1]$ ordered by the order relation $\leq_{\mathcal{P}}$ given by Def.~\ref{def_order_power}.

\noindent Set $[a_0,b_0]:=\left[0,\frac{1}{4}\right]$ and $[a_1,b_1]:=\left[\frac{3}{4},1\right]$. 
Let us define
\begin{align*}
L_0&:=\{[a_0,b_0],[a_1,b_1]\}=\left \{\left[0,\frac{1}{4}\right], \left[\frac{3}{4},1\right]\right \}\\
\\
L_1&:=\left\{[a_0,b_0], \left[a_2, b_2\right], [a_1,b_1]\right\}=L_0\cup\{[a_2,b_2]\},\\
\end{align*}
where $[a_2,b_2]$ is intermediate between $[a_0,b_0]$ and $[a_1,b_1]$. Notice that $[a_0,b_0] <_\mathcal{P} [a_2,b_2]<_\mathcal{P} [a_1,b_1]$.
Let us proceed by setting
\[
L_2:=\left\{[a_0,b_0],[a_3, b_3], [a_2,b_2],[a_4,b_4],[a_1,b_1]\right\}=L_1\cup\{[a_3,b_3],[a_4,b_4]\}
\]
where 
\begin{itemize}
\item $[a_3,b_3]$ is intermediate between $[a_0,b_0]$ and $[a_2,b_2]$, 
\item $[a_4,b_4]$ is intermediate between $[a_2,b_2]$ and $[a_1,b_1]$.
\end{itemize}
So, we see that $[a_0,b_0]<_\mathcal{P}[a_3,b_3]  <_\mathcal{P} [a_2,b_2]<_\mathcal{P}[a_4,b_4]<_\mathcal{P} [a_1,b_1]$.

\noindent Proceeding inductively, for $n\ge 2$, given $L_n=\{I_0, I_1,\cdots , I_{2^n}\}$ with $I_0<_\mathcal{P} I_1<_\mathcal{P}\cdots <_\mathcal{P} I_{2^n}$, we set
\[
L_{n+1}:=L_n\cup\{I_{2^n+1}, I_{2^n+2},\cdots , I_{2^{n+1}}\}
\]
where $I_{2^n+k}=[a_{2^n+k}, b_{2^n+k}]$ is intermediate between $I_{k-1}$ and $I_{k}$ for every $k\in\{1,\ldots,2^n\}$.
So, it holds that $I_{k-1} <_\mathcal{P} I_{2^n+k} <_\mathcal{P} I_k$ for every $k\in\{1,\ldots,2^n\}$.
Notice that $\left|L_n\right|=2^n+1$ for every $n\geq 0$, and that $L_n\subsetneq L_m$ whenever $n<m$.
Set 
$$
A_n:=[0,1]\setminus \cup L_n.
$$ 
Observe that $A_n$ is an open set and $A_{n+1}\subseteq A_{n}$ for all $n\in\mathbb{N}_0$. Moreover, we can choose the intervals $\{[a_n,b_n]\}_{n>1}$ so that 
\begin{equation}\label{eq_mis_compl}
    \lim_{n\to\infty} m(A_n)=0,
\end{equation}
where $m$ indicates the Lebesgue measure on $[0,1]$.

\noindent Finally, we set 
\[
L:=\bigcup_{n\geq 0}L_n=\{[a_n,b_n]\}_{n\in\mathbb{N}_0}.
\]
Let us now define $f:[0,1]\to [0,1]$ as:
\begin{equation}\label{dis_func}
    f(x)=
    \begin{cases}
        a_n &\quad\text{if } x\in[a_n,b_n] \text{ for some }n\in\mathbb{N}_0\\
        0 &\quad\text{if } x\in[0,1]\setminus \cup L
    \end{cases}
\end{equation}
From the construction of $L$ it is clear that $(L,\le_\mathcal{P})$ is linearly ordered and for every $I, J\in L$ such that $I<_\mathcal{P} J$, there is $K\in L$ such that
\[
I<_\mathcal{P} K<_\mathcal{P} J,
\]
which means that $(L,<_\mathcal{P})$ is also densely ordered.
Since $[a_0,b_0]<_\mathcal{P} [a_n,b_n]<_\mathcal{P} [a_1,b_1]$ for every $n\ge2$, we have that $(L,\le_\mathcal{P})$ is order isomorphic to $([0,1]\cap \mathbb{Q}, \leq)$.

\noindent It remains to show that $(\mathfrak{C}_f,\preceq)$ is order isomorphic to $(L,\le_\mathcal{P})$.
Notice that the points $a_n$ for every $n\in\mathbb{N}_0$ are the only fixed points for the map $f$.
We want to prove that
$$
\mathfrak{C}_f=
\{\{a_n\}: n\in\mathbb{N}_0\}.
$$
\noindent Let us show that for every $[a_p,b_p],[a_q,b_q]\in L$ with $[a_p,b_p]<_\mathcal{P} [a_q,b_q]$ it holds that $[a_p]\prec [a_q]$.
\begin{enumerate}
\item Let us first prove that $[a_p]\preceq [a_q]$, that is, $a_q\,\mathcal{C}\, a_p$.   

Fix $\varepsilon>0$.
If $\varepsilon> a_q-b_p$ then the three points $a_q, b_p, a_p$ form an $\varepsilon$-chain from $a_q$ to $a_p$.
Indeed, $d(f(a_q),b_p)=d(a_q,b_p)<\varepsilon$ and $f(b_p)=a_p$.
So suppose that $\varepsilon<a_q-b_p$. 
By \eqref{eq_mis_compl} and since $(L,<_\mathcal{P})$ is densely ordered, there is $n$ sufficiently large that we can choose $n-2$ intervals 
\[
J_{n-2}<_\mathcal{P} J_{n-1}<_\mathcal{P}\cdots <_\mathcal{P} J_1
\]
in $L$ with $J_1=[x_1,y_1]$, $J_2=[x_2,y_2]$, $\ldots$, $J_{n-2}=[x_{n-2},y_{n-2}]$, such that
\begin{align*}
    &d(a_q,y_1)<\varepsilon\\
    &d(x_1,y_2)<\varepsilon\\
    &\quad\quad\vdots\\
    &d(x_{n-3},y_{n-2})<\varepsilon\\
    &d(x_{n-2},b_p)<\varepsilon
\end{align*}
From the definition of $f$ given by \eqref{dis_func}, we have that $f(a_q)=a_q$ and $f(y_i)=x_i$ for every $i=1,\ldots,n-2$. Therefore, we can rewrite the previous inequalities as
\begin{align*}
    &d(f(a_q),y_1)<\varepsilon\\
    &d(f(y_1),y_2)<\varepsilon\\
    &\quad\qquad\vdots\\
    &d(f(y_{n-3}),y_{n-2})<\varepsilon\\
    &d(f(y_{n-2}),b_p)<\varepsilon
\end{align*}
Since $f(b_p)=a_p$, it follows that the points
$$
a_q,y_1,y_2,\ldots,y_{n-2},b_p, a_p
$$
form an $\varepsilon$-chain from $a_q$ to $a_p$. 
We conclude that $[a_p]\preceq [a_q]$.

\item It remains to show that $[a_q]\not\preceq [a_p]$.
Notice that, by \eqref{dis_func} we have $f(x)\le x$ for every $x\in [0,1]$.
Fix $\varepsilon>0$ such that $\varepsilon< b_p-a_p\leq\frac{1}{4}$.
Then, there can be no $\varepsilon$-chain 
\[
a_p=x_0,x_1,x_2,\ldots,x_n=a_q
\]
from $a_p$ to $a_q$.
In fact, suppose by contradiction that there is an $\varepsilon$-chain  $x_0,\ldots,x_n$ from $a_p$ to $a_q$. If $x_1\ge x_0=a_p$, since 
\begin{equation}\label{d(f(x_0),x_1)<epsilon}
d(f(x_0),x_1)<\varepsilon,
\end{equation}
then $x_1\in [a_p,b_p]$. Indeed, since $f(x_0)=f(a_p)=a_p$, inequality (\ref{d(f(x_0),x_1)<epsilon}) becomes $d(a_p,x_1)<\varepsilon$.
And again, if $x_2\ge x_0$, the condition 
$$
d(f(x_1),x_2)=d(a_p,x_2)<\varepsilon
$$
implies that $x_2\in [a_p,b_p]$.
So, let $1\le j\le n-1$ be the first index for which 
$$
x_j< f(x_{j-1})=a_p=x_0.
$$
If $x_j\in [0,1]\setminus \cup L$, then $f(x_j)=0$ and since $\varepsilon<\frac 1 4$, it follows that 
$$
f(x_{k})=0 \quad \forall\, k=j,\ldots,n-1.
$$ 
Hence $a_q\in[0,\frac 1 4]$, which is impossible.
If $x_j\in L$, since $f(x_j)\le  x_j< a_p$, then $x_k < b_p$ for every $j<k\le n$. Therefore, $d(f(x_{n-1}),a_q)\ge d(a_p,a_q)>\varepsilon$ which is a contradiction.
\end{enumerate}

\noindent In particular, from the definition of the map $f$ given by \eqref{dis_func} and by points 1. and 2. we have that
$
\mathfrak{C}_f=
\{\{a_n\}:n\in\mathbb{N}_0\}
$.

\noindent Finally, by points 1. and 2. it follows that the application
$$
[a_n,b_n]\mapsto \{a_n\} \qquad (n\in\mathbb{N}_0)
$$
is an order isomorphism between $(L,\le_\mathcal{P})$ and $(\mathfrak{C}_f,\preceq)$.
\end{proof}

\end{document}
